\section{Properties of the proportional electability cost $\Cfun{v}{\delta}$}

Fix $\delta>0$. Recall $\Cfun{v}{\delta}\equiv \PhiStd(v)/\PhiStd(v-\delta)$ (Definition~\ref{def:Cbasic}).

The object $\Cfun{v}{\delta}$ isolates the \emph{general-election} penalty from taking action $a=1$:
it is the proportional decrease in general-election win probability when the nominee's index shifts
from $v$ to $v-\delta$. Its properties are useful for establishing monotonicity and for characterizing
best responses via $\Gfun{v}{\Sagg}\ge \Cfun{v}{\delta}$.


\begin{lemma}[Properties of the proportional electability cost]\label{lem:Cbasic}
Fix $\delta > 0$. The proportional electability cost, given by
\[
    \Cfun{v}{\delta} = \frac{\PhiStd(v)}{\PhiStd(v-\delta)},
\]
satisfies the following properties:
\begin{enumerate}
    \item $\Cfun{v}{\delta} > 1$;
    \item $\lim_{v\to+\infty}\Cfun{v}{\delta} = 1$ and $\lim_{v\to -\infty} \Cfun{v}{\delta} = +\infty$;
    \item $\Cfun{v}{\delta}$ is strictly increasing in $\delta$ and strictly decreasing in $v$.
\end{enumerate}
\end{lemma}


Part (ii) formalizes that electability costs are negligible for very high-valence candidates
(both $\PhiStd(v)$ and $\PhiStd(v-\delta)$ are near $1$), but potentially very large for low-valence
candidates (where $\PhiStd(v-\delta)$ can be extremely small). Part (iii) ensures the GE penalty from
$a=1$ becomes more severe as $\delta$ increases and as $v$ falls.


\begin{proof}
Fix $\delta>0$.

\begin{enumerate}
    \item Since $v-\delta < v$ and $\PhiStd(\cdot)$ is strictly increasing, we have
    $\PhiStd(v-\delta) < \PhiStd(v)$. Moreover $\PhiStd(v-\delta)>0$ for all finite $v$.
    Hence
    \[
        \Cfun{v}{\delta}=\frac{\PhiStd(v)}{\PhiStd(v-\delta)} > 1.
    \]

    \item As $v\to +\infty$, $\PhiStd(v)\to 1$ and $\PhiStd(v-\delta)\to 1$, so
    \[
        \lim_{v\to +\infty}\Cfun{v}{\delta}= \frac{1}{1}=1.
    \]
    As $v\to -\infty$, $\PhiStd(v)\to 0$ and $\PhiStd(v-\delta)\to 0$. Consider
    \[
        \Cfun{v}{\delta}=\frac{\PhiStd(v)}{\PhiStd(v-\delta)}
        =\frac{1/\PhiStd(v-\delta)}{1/\PhiStd(v)}.
    \]
    Using the standard normal left-tail asymptotic $\PhiStd(x)\sim \frac{\phiStd(x)}{-x}$ as
    $x\to -\infty$, we obtain
    \[
        \frac{\PhiStd(v)}{\PhiStd(v-\delta)}
        \sim
        \frac{\phiStd(v)}{\phiStd(v-\delta)}\cdot \frac{-(v-\delta)}{-v}
        =
        \exp\!\left(\delta v-\frac{\delta^2}{2}\right)\cdot\left(1-\frac{\delta}{v}\right),
        \qquad v\to -\infty,
    \]
    which diverges to $+\infty$ because $\delta v\to -\infty$ and the exponential term dominates.
    Therefore $\lim_{v\to -\infty}\Cfun{v}{\delta}=+\infty$.

    \item (Monotonicity in $\delta$.) For fixed $v$, the denominator $\PhiStd(v-\delta)$ is strictly
    decreasing in $\delta$ (since $v-\delta$ falls as $\delta$ rises), while the numerator does not
    depend on $\delta$. Since the denominator is positive, $\Cfun{v}{\delta}$ is strictly increasing
    in $\delta$.

    (Monotonicity in $v$.) Differentiate $\log \Cfun{v}{\delta}$ with respect to $v$:
    \[
        \frac{\partial}{\partial v}\log \Cfun{v}{\delta}
        = \frac{\phiStd(v)}{\PhiStd(v)} - \frac{\phiStd(v-\delta)}{\PhiStd(v-\delta)}.
    \]
    The function $m(x)\equiv \frac{\phiStd(x)}{\PhiStd(x)}$ (the inverse Mills ratio) is strictly
    decreasing in $x$, so $m(v) < m(v-\delta)$ for $\delta>0$. Hence
    \[
        \frac{\partial}{\partial v}\log \Cfun{v}{\delta} < 0,
    \]
    which implies $\frac{\partial}{\partial v}\Cfun{v}{\delta}<0$ and therefore $\Cfun{v}{\delta}$
    is strictly decreasing in $v$.
\end{enumerate}
\end{proof}